\chapter{INTRODUCTION}
\label{ch:introduction}

Cryptography is a foundation of modern computing systems. It protects communication, supports authentication, preserves data integrity, and enables privacy-preserving applications. From secure messaging and online banking to blockchains and verifiable computation, software systems increasingly rely on cryptographic mechanisms to establish trust when data, computation, and infrastructure may be exposed to adversarial behavior. Cryptographic mechanisms are widely used to provide core security services such as confidentiality, integrity, authentication, and non-repudiation \cite{barker2016guideline}. These mechanisms are often built on rigorous mathematical foundations, but their guarantees do not automatically transfer to deployed software. They depend on the accuracy of the implementation, the assumptions made by developers, and the ability of testing and analysis tools to expose mistakes before deployment.

Although cryptography is primarily developed to protect sensitive information and secure communications, it can also be exploited for malicious purposes. In ransomware attacks, adversaries use cryptographic functions to encrypt victims' files and demand payment in exchange for restoring access \cite{beaman2021ransomware}. This dual use creates a significant challenge for cybersecurity: the same mechanisms that strengthen privacy and system security can also conceal malicious behavior. The use of cryptographic operations within malware can obscure critical functionality, complicate reverse engineering, and hinder security analysts' efforts to understand the attack and develop effective detection and mitigation strategies.

The gap between cryptographic design and real-world implementation creates a persistent source of risk. A protocol may be secure in theory, while its implementation may still contain errors that undermine the intended guarantee. These errors can arise from incorrect arithmetic, missing validation checks, inconsistent data conversion, incomplete constraint generation, or unexpected interactions between different software components. Studies of real-world cryptographic software show that implementation mistakes, API misuse, and incorrect validation logic can compromise applications even when they rely on standard cryptographic mechanisms \cite{georgiev2012most,lazar2014does}. Unlike conventional software bugs, cryptographic implementation errors may not produce crashes or visibly incorrect outputs. A program may execute normally, generate plausible results, and still fail to enforce the property it was designed to provide. This makes implementation-level analysis essential for understanding whether cryptographic software is secure in practice.

Analyzing cryptographic software is also difficult because its structure differs from many ordinary programs. Implementations often contain dense mathematical computation, optimized arithmetic routines, carefully formatted inputs, and dependencies that are not immediately visible from source code or execution traces. In many security settings, analysts must work with compiled binaries rather than clean source code. Malware analysis, software auditing, reverse engineering, and deployment validation frequently require reasoning about executables affected by compiler choices, optimization levels, architecture differences, and library dependencies. These factors make it difficult to rely only on protocol descriptions or source-level inspection \cite{shoshitaishvili2016sok}. A practical security analysis must connect the intended cryptographic behavior with the behavior of the software that is actually compiled, deployed, and executed.

This challenge becomes even more pronounced for modern cryptographic systems such as zero-knowledge proofs. Zero-knowledge proofs allow a prover to convince a verifier that a statement is true without revealing additional information beyond the truth of the statement. They have become a central tool for privacy-preserving and verifiable computation, with applications in blockchains, authentication, confidential computation, and other settings where correctness or privacy must be established without revealing sensitive information. However, implementations of these systems combine several complex layers, including finite-field arithmetic, constraint systems, witness generation, proving procedures, verification logic, and serialization formats. A mistake in any layer can create a mismatch between the statement the developer intended to prove and the statement actually enforced by the implementation.

General-purpose testing techniques are not sufficient for all of these problems. Software testing and automated analysis have long been used to expose implementation defects, and fuzzing has become an effective approach for finding software bugs by generating and mutating inputs \cite{manes2019art}. However, cryptographic applications often require structured and interdependent inputs. This is especially true for zero-knowledge proof programs, where circuits, witnesses, proving keys, verification keys, and proofs must remain consistent with one another. Random mutation can easily break these relationships before execution reaches meaningful cryptographic logic. Therefore, testing zero-knowledge proof implementations requires techniques that combine program feedback with domain knowledge about the cryptographic workflow.

These observations motivate systematic and automated techniques for improving confidence in cryptographic implementations. The work presented here addresses three connected problems. The first is how to evaluate and improve understanding of cryptographic function identification in binary programs. The second is how to detect and locate logic errors in zkSNARK implementations. The third is how to test zero-knowledge proof binary applications when source-level visibility is limited, in order to support the security analysis of deployed systems. Together, these directions study cryptographic implementation security across several levels of abstraction, from binary-level detection to protocol-aware analysis for zero-knowledge proof protocols.

\autoref{ch:cryptobinary} studies cryptographic function identification in binaries. Cryptographic routines appear in secure applications, libraries, malware, and ransomware, so detecting them is valuable for reverse engineering, forensic analysis, auditing, and security evaluation. Existing tools rely on signatures, static analysis, dynamic traces, symbolic reasoning, data-flow features, or machine-learning models, but their evaluations are often difficult to compare. This chapter examines available tools, categorizes their detection strategies, and performs reproduction and replication studies using a newly developed benchmarking framework. The evaluations consider different compilers, optimization levels, obfuscation strategies, and algorithm variants to reveal where current approaches succeed, where they fail, and what gaps remain for future work.

\autoref{ch:snarkprobe} turns to the correctness of zero-knowledge proof implementations, with a focus on zkSNARK systems. zkSNARK implementations must correctly connect high-level statements to low-level constraints, witness generation, finite-field arithmetic, proving, and verification. A subtle mistake in any of these components may cause the implementation to enforce a different computation from the one intended. This chapter introduces SNARKProbe, an automated security analysis framework that combines constraint checking with fuzzing-based testing to detect and locate cryptographic logic errors. The framework helps determine whether the implemented constraint system matches the intended computation and where an inconsistency originates when the behavior is incorrect.

\autoref{ch:greybox} investigates grey-box differential fuzzing for zero-knowledge proof binary applications. Many cryptographic applications, including zero-knowledge proof systems, are distributed as closed-source software, limiting access to source-level information and making conventional code inspection difficult. This chapter develops a grey-box approach that uses structured input generation, coverage monitoring, and error localization to guide testing toward security-relevant code. Differential testing is then used to expose inconsistencies across related executions or components and to help localize faults in closed-source \zksnark applications.

Taken together, these works approach cryptographic implementation security from binary analysis, automated checking, and fuzzing. The first contribution clarifies the current state of binary-level cryptographic function detection. The second provides a framework for finding logic errors in zkSNARK implementations. The third extends testing to deployed zero-knowledge proof binaries through grey-box differential analysis. By combining these perspectives, this work aims to strengthen the practical security of cryptographic software and support future tools that can more effectively analyze, test, and secure real-world implementations.